\documentclass[11pt,a4paper]{article}
\usepackage[utf8]{inputenc}
\usepackage[T1]{fontenc}
\usepackage{lmodern}
\usepackage[english]{babel}
\usepackage{amsmath,amssymb}
\usepackage[margin=2.4cm]{geometry}
\usepackage{booktabs}
\usepackage{hyperref}

\newcommand{\R}{\mathbb{R}}
\newcommand{\E}{\mathbb{E}}
\newcommand{\norm}[1]{\left\lVert #1 \right\rVert}
\newcommand{\fonte}[1]{\par\noindent\textit{Source: #1}\par\smallskip}
\newcommand{\esterna}[1]{\par\noindent\textit{External source (not in the course materials): #1}\par\smallskip}

\title{Double descent in linear models\\\large Theory for Assignment 4 (SMML 2025/26, Prof.\ Cesa-Bianchi)}
\author{Filippo Ferroni}
\date{Version 4, 25 October 2026}

\begin{document}
\maketitle

\paragraph{How to use this document.}
Reading time is about 40 minutes. Sections 2--6 are the minimum theory to be able to explain at the oral exam.
Section 7 links the theory to the gradient descent extension, Section 8 describes what we expect to see in the plots,
Section 9 the experimental design. At the end: likely oral questions with answers, and solutions to the check questions.
Every claim states its source; when a concept is \textbf{not} in the uploaded course materials (Cesa-Bianchi's notes 1--15) it is
flagged explicitly as an external source.

\tableofcontents

%--------------------------------------------------------------------
\section{Notation and correspondences}

\begin{center}
\begin{tabular}{@{}llll@{}}
\toprule
Concept & Course (slides 9) & Assignment / paper & Code \\
\midrule
number of training examples & $m$ & $n$ & \texttt{m} \\
number of features (complexity) & $d$ & $d$ (assignment) / $N$ (paper) & \texttt{d} \\
design matrix & $X \in \R^{m\times d}$ & & \texttt{X} \\
labels & $y \in \R^m$ & & \texttt{y} \\
ridge parameter & $\alpha$ & $\lambda$ (other texts) & \texttt{alpha} \\
interpolation threshold & $d = m$ & $N = n$ & \\
\bottomrule
\end{tabular}
\end{center}

\textbf{Careful:} in the paper $d$ is the dimension of the \emph{input} $x$, while the number of features/parameters is $N$ and the number of samples is $n$ (p.~2 and p.~4).

In the report we state once: ``we follow the course notation; the number of examples $n$ in the assignment corresponds
to $m$''. The rows of $X$ are the examples: $X^\top = [x_1,\dots,x_m]$ (slides 9, p.~5).
A linear predictor for regression is $h(x) = w^\top x$ with $w \in \R^d$ (slides 9, p.~1 and p.~4).
The intercept is handled by adding a constant feature equal to $1$ (slides 9, p.~1) or by centering the data.

%--------------------------------------------------------------------
\section{Statistical setup: risk, training error, test error}
\fonte{slides ``4. statRisk'' (p.~1--4) and ``1. intro'' (p.~3--5).}

\begin{itemize}
  \item Every example $(x,y)$ is an independent draw from a fixed but unknown distribution $\mathcal{D}$ (statRisk, p.~1).
        Training and test sets are therefore statistical samples.
  \item \textbf{Statistical risk} of a predictor $h$ with respect to the loss $\ell$ (statRisk, p.~1):
        $\ell_{\mathcal{D}}(h) = \E\bigl[\ell(h(X), Y)\bigr]$. For us $\ell(\hat y, y) = (\hat y - y)^2$ and $h(x) = w^\top x$.
  \item \textbf{Bayes optimal predictor} $f^\ast$: the best possible one given $\mathcal{D}$. For the square loss
        $f^\ast(x) = \E[Y \mid X = x]$ and the \textbf{Bayes risk} is the expected conditional variance of the label
        (statRisk, p.~1--2). In our synthetic dataset, with $y = w^{\star\top}x + \varepsilon$ and noise variance $\sigma^2$,
        the Bayes risk is exactly $\sigma^2$: \textbf{no model can have expected test error below $\sigma^2$}.
        We will use it as a reference line in the plots.
  \item \textbf{Training error}: $\ell_S(h) = \frac{1}{m}\sum_{t=1}^{m} \ell(h(x_t), y_t)$; for us
        $\ell_S(w) = \frac{1}{m}\norm{Xw - y}^2$. ERM picks the predictor minimizing it (intro, p.~4).
  \item \textbf{Test error}: the same average on an independent test set $S'$; it estimates the risk $\ell_{\mathcal{D}}(h)$,
        and by Chernoff--Hoeffding the estimate is accurate when the test set is large (statRisk, p.~3--4).
        The test set is never used for training or for choosing hyperparameters.
\end{itemize}

\paragraph{Overfitting and underfitting in the course.}
Intro, p.~5: if the algorithm returns a predictor with high training error it is \emph{underfitting}; with low training error
but high test error it is \emph{overfitting}. The course notes that overfitting is expected when the predictor class is too
rich compared to $m$ (``$\log_2|F| \gg m$'', intro p.~5). \textbf{Double descent is precisely a case where this classical
intuition is not enough}: for $d \gg m$ the training error is $0$ but the test error goes down again.

%--------------------------------------------------------------------
\section{Least squares with $d \le m$: the formula in the slides}
\fonte{slides 9, p.~4--5.}

The linear regression predictor is the ERM with respect to the square loss:
\[
  w_S = \operatorname*{argmin}_{w \in \R^d} \norm{Xw - y}^2 .
\]
$F(w) = \norm{Xw - y}^2$ is convex, so the minimum is where the gradient vanishes:
\[
  \nabla F(w) = 2X^\top(Xw - y) = 0
  \quad\Longleftrightarrow\quad
  X^\top X\, w = X^\top y \qquad \text{(normal equations)}.
\]
If $X^\top X$ is invertible:
\[
  \boxed{w_S = (X^\top X)^{-1} X^\top y}
\]
\textbf{When is it invertible?} The slides say: if and only if $x_1,\dots,x_m$ span $\R^d$.
At least $d$ vectors are needed to span $\R^d$, so we need $m \ge d$.

\paragraph{Geometric interpretation.}
The vector of training predictions is $Xw_S = X(X^\top X)^{-1}X^\top y = Py$, where $P$ is the orthogonal projection
onto the column space of $X$. The more features we add, the larger that space and the closer $Py$ gets to $y$:
the training error \textbf{cannot increase} as $d$ grows.

\paragraph{The case $d = m$.}
If $X$ is square and invertible, $w_S = X^{-1}y$ and $Xw_S = y$: training error exactly $0$.
The model already \textbf{interpolates} at $d = m$. Hence the name ``interpolation threshold''.

%--------------------------------------------------------------------
\section{Least squares with $d > m$: the minimum-norm solution}
\esterna{Moore--Penrose pseudo-inverse, standard linear algebra textbooks (e.g.\ Strang, \emph{Introduction to Linear
Algebra}, chapter on the pseudo-inverse and SVD); use of the minimum-norm solution: Belkin et al.\ 2019, Sec.~2 (p.~4: ``when the minimizer is not unique \dots{} we choose the minimizer whose coefficients have the minimum $\ell_2$ norm'') and p.~10 (``minimum norm linear regression in the feature space'').}

\subsection{Why the slides' formula no longer works}
$X^\top X$ is a $d \times d$ matrix, but its rank equals the rank of $X$, which is at most $\min(m,d) = m$.
If $d > m$ then $\operatorname{rank}(X^\top X) \le m < d$: the matrix is \textbf{singular} and $(X^\top X)^{-1}$ does not exist.
Equivalently: $m$ vectors cannot span $\R^d$ when $m < d$.

\subsection{Infinitely many interpolating solutions}
With $d > m$ (and linearly independent rows of $X$, the typical case with continuous data), the system $Xw = y$ has
$m$ equations and $d$ unknowns: it has \textbf{infinitely many} exact solutions. If $w_0$ is a solution, so is $w_0 + z$
for every $z$ in the null space of $X$ ($Xz = 0$), which has dimension $d - m$. All of them have zero training error.

\subsection{Which one to pick: the minimum-norm solution}
\[
  \boxed{w^{+} = \operatorname*{argmin}\bigl\{\norm{w} : Xw = y\bigr\} = X^\top (X X^\top)^{-1} y}
\]
\textbf{Check that it interpolates:} $Xw^{+} = XX^\top(XX^\top)^{-1}y = y$. ($XX^\top$ is $m\times m$ and invertible when
the rows of $X$ are independent.)

\textbf{Why it has minimum norm (short proof, know how to do it).}
$w^{+} = X^\top a$ with $a = (XX^\top)^{-1}y$, so $w^{+}$ is a linear combination of the rows of $X$.
Any other solution is $w = w^{+} + z$ with $Xz = 0$, i.e.\ $z$ is orthogonal to every row of $X$, hence to $w^{+}$.
By Pythagoras:
\[
  \norm{w}^2 = \norm{w^{+}}^2 + \norm{z}^2 \ge \norm{w^{+}}^2 .
\]

\subsection{One formula for everything: the pseudo-inverse}
The pseudo-inverse $X^{+}$ is computed from the SVD $X = U\Sigma V^\top$: $X^{+} = V\Sigma^{+}U^\top$, where $\Sigma^{+}$
inverts the non-zero singular values and leaves the others at zero. Then:
\begin{itemize}
  \item if $d \le m$ and $X$ has full rank: $X^{+}y = (X^\top X)^{-1}X^\top y = w_S$ (the slides' formula);
  \item if $d > m$ and $X$ has rank $m$: $X^{+}y = X^\top(XX^\top)^{-1}y = w^{+}$.
\end{itemize}
So in numpy \texttt{w = np.linalg.pinv(X) @ y} works for \textbf{every} $d$.
It is linear algebra, not a library ``fit'', so it respects the ``from scratch'' rule (assignment FAQ, p.~16).

\paragraph{Numerical notes (important for the code).}
\begin{itemize}
  \item Do not use \texttt{np.linalg.inv(X.T @ X)}: forming $X^\top X$ squares the condition number, and near $d=m$
        the results become numerically unreliable.
  \item \texttt{np.linalg.pinv} only zeroes singular values below a tiny threshold (\texttt{rcond}, about $10^{-15}$
        relative to the largest). Near $d=m$ the small singular values are \emph{not} zeroed: \texttt{pinv} does not
        regularize, so the peak stays visible. This is the behaviour we want to show.
\end{itemize}

\subsection{Minimal numerical example}
$m=1$, $d=2$: a single example $x = (1,1)$, $y = 2$. The solutions of $w_1 + w_2 = 2$ are infinitely many: $(2,0)$,
$(0,2)$, $(1,1)$, $(3,-1)$\dots{} The minimum-norm solution is $w^{+} = x^\top(xx^\top)^{-1}y = (1,1)\cdot\frac{2}{2} = (1,1)$,
with norm $\sqrt{2}$, versus norm $2$ for $(2,0)$. It spreads the weight evenly instead of ``betting everything'' on
one feature.

%--------------------------------------------------------------------
\section{Why the test error blows up at $d \approx m$}
\fonte{slides 9, p.~5: ``when $X^\top X$ is nearly singular, $w_S$ is highly sensitive to perturbations of the training
set. This instability increases the estimation error (or variance)''. This is the hook to the course.}

\subsection{Intuition with a $2\times 2$ example}
$m = d = 2$, rows $x_1 = (1,0)$, $x_2 = (1,\varepsilon)$ with small $\varepsilon$ (almost parallel rows, $X$ nearly
singular). Labels $y = (1,\, 1+\delta)$, where $\delta$ is a small noise.
Solving $Xw = y$: $w_1 = 1$ and $w_1 + \varepsilon w_2 = 1+\delta$, so
\[ w_2 = \delta/\varepsilon . \]
With $\varepsilon = 0.01$ and noise $\delta = 0.1$ we get $w_2 = 10$: a small noise produces a huge weight.
On a new point, e.g.\ $x = (0,1)$, the prediction is $10$ instead of roughly $0$. That is exactly the peak.

\subsection{General version with eigenvalues}
Let $\lambda_1 \ge \dots \ge \lambda_d \ge 0$ be the eigenvalues of $X^\top X$ and $v_1,\dots,v_d$ the (orthonormal)
eigenvectors. The minimum-norm solution reads
\[
  w^{+} = \sum_{i:\,\lambda_i > 0} \frac{v_i^\top X^\top y}{\lambda_i}\, v_i .
\]
Assume noisy linear labels: $y = Xw^\star + \varepsilon$. The noise part is
$\sum_i \frac{v_i^\top X^\top \varepsilon}{\lambda_i} v_i$; each term has norm
$\frac{|u_i^\top \varepsilon|}{\sqrt{\lambda_i}}$, with $u_i = Xv_i/\sqrt{\lambda_i}$ a unit vector.
\textbf{The noise is divided by $\sqrt{\lambda_i}$.}
\begin{itemize}
  \item $d \ll m$: $X$ is ``tall and thin'', all eigenvalues are large (about $m$ for standardized features):
        the noise is damped.
  \item $d \approx m$: $X$ is nearly square and random, the smallest eigenvalue is very close to $0$:
        the noise is amplified, $\norm{w}$ blows up, and so does the test error.
  \item $d \gg m$: $X$ is ``wide'', the $m$ non-zero eigenvalues move away from $0$ again and the noise is spread over
        many directions: $\norm{w^{+}}$ and the test error decrease. This is the \textbf{second descent}.
\end{itemize}
\esterna{the behaviour of the smallest eigenvalue of a random matrix as a function of $d/m$ is a random matrix theory
result (Marchenko--Pastur law). We do not prove it: we show it empirically by plotting $\norm{w}$ against $d$, which is
easier to defend at the oral.}

\subsection{An exact formula for the classical regime (optional)}
\esterna{Hastie, Montanari, Rosset, Tibshirani, ``Surprises in high-dimensional ridgeless least squares interpolation'',
Annals of Statistics 2022.}
If $x \sim \mathcal{N}(0, I_d)$, $y = w^{\star\top}x + \varepsilon$ with noise variance $\sigma^2$ and $d < m-1$, then
\[
  \E\bigl[\norm{w_S - w^\star}^2\bigr] = \sigma^2\,\frac{d}{m - d - 1},
\]
which \textbf{diverges as $d \to m-1$}. No need to prove it; it is useful to say at the oral that the peak is not a
numerical artefact but a statistical phenomenon predicted by theory.

%--------------------------------------------------------------------
\section{Ridge regression}
\fonte{slides 9, p.~5.}

\subsection{Formula}
\[
  w_{S,\alpha} = \operatorname*{argmin}_{w \in \R^d} \norm{Xw - y}^2 + \alpha \norm{w}^2 ,
  \qquad \alpha > 0 .
\]
Gradient: $2X^\top(Xw - y) + 2\alpha w = 0 \;\Rightarrow\; (X^\top X + \alpha I)\,w = X^\top y$, hence
\[
  \boxed{w_{S,\alpha} = (\alpha I + X^\top X)^{-1} X^\top y}
\]

\subsection{Why it removes the peak}
The eigenvalues of $\alpha I + X^\top X$ are $\alpha + \lambda_i \ge \alpha > 0$ (slides 9, p.~5): the matrix is always
invertible, for every $d$, including $d > m$. In the eigenvector basis:
\[
  w_{S,\alpha} = \sum_{i} \frac{v_i^\top X^\top y}{\lambda_i + \alpha}\, v_i .
\]
We divide by $\lambda_i + \alpha$ instead of $\lambda_i$: even when $\lambda_i \approx 0$ the denominator stays at least
$\alpha$, so the noise cannot be amplified without bound. The peak is damped or disappears.

\subsection{The price: bias}
\begin{itemize}
  \item $\alpha \to 0$: we recover the minimum-norm solution $w^{+}$ (also for $d > m$; along directions with
        $\lambda_i = 0$ the numerator $v_i^\top X^\top y$ is $0$). \emph{Standard fact, not in the slides.}
  \item $\alpha \to \infty$: $w_{S,\alpha} \to 0$, the predictor ignores the data (slides 9, p.~5: $\alpha$ controls the bias).
  \item In between: $\alpha$ trades variance (instability) for bias (underfitting). With too large an $\alpha$ the test
        error goes up \emph{for every} $d$.
\end{itemize}
\fonte{Belkin et al.\ 2019, p.~10: regularization ``can both prevent interpolation and change the effective capacity of the function class, thus attenuating or masking the interpolation peak''.}
\esterna{Nakkiran, Venkat, Kakade, Ma, ``Optimal regularization can mitigate double descent'', ICLR 2021: with the optimal
$\alpha$ the test error curve becomes monotone.}

\subsection{Computing it in practice}
\begin{itemize}
  \item Primal: \texttt{np.linalg.solve(X.T @ X + alpha * np.eye(d), X.T @ y)} ($d \times d$ system).
  \item Dual, cheaper for $d > m$: $w_{S,\alpha} = X^\top(XX^\top + \alpha I)^{-1}y$ ($m\times m$ system).
        The two forms agree because $X^\top(XX^\top + \alpha I) = (X^\top X + \alpha I)X^\top$.
  \item \texttt{solve} is preferable to computing the explicit inverse (more stable and faster).
\end{itemize}

\subsection{Choosing $\alpha$ without data leakage}
\begin{itemize}
  \item For the main plot: a few curves with fixed $\alpha$ (e.g.\ $10^{-3}, 10^{-1}, 10$) next to the minimum-norm curve.
        No data-driven choice, no leakage.
  \item Time permitting: $\alpha$ chosen for each $d$ on a validation split (or $K$-fold cross-validation) taken
        \textbf{only} from the training set (slides ``6. crossVal'', p.~2--3: tuning on a validation set; p.~3--4: nested CV),
        as in the professor's \texttt{CrossValidation.ipynb} notebook.
        The test set is never used for selection.
\end{itemize}

%--------------------------------------------------------------------
\section{Link with gradient descent (extension)}
\fonte{slides 10 (online gradient descent) for the algorithm; the minimum-norm result is external (see below).}

Gradient descent on $F(w) = \norm{Xw - y}^2$ started at $w_1 = 0$:
\[
  w_{k+1} = w_k - \eta\, 2X^\top(Xw_k - y).
\]
\begin{itemize}
  \item Every update is $X^\top(\dots)$, a linear combination of the rows of $X$. Starting from $0$, $w_k$ always stays
        in the row space.
  \item With a small enough step size ($\eta < 1/\lambda_1$), for $d > m$ GD converges to a solution that interpolates
        \emph{and} lies in the row space: by the proof in Section 4.3, this is exactly $w^{+}$.
  \item Stopping before convergence (early stopping) has an effect similar to ridge: directions with small $\lambda_i$
        are learned last, so the noise along them is not amplified.
\end{itemize}
\fonte{implicit bias of GD towards minimum norm: Belkin et al.\ 2019, p.~11 (``SGD initialized at zero also converges to the minimum norm solution'').}
The extension is done \textbf{only} once the report and repo are complete.

%--------------------------------------------------------------------
\section{What we expect in the plots}

\subsection{The classical U-curve and double descent}
\begin{itemize}
  \item \textbf{Classical regime} ($d < m$): the course decomposes the risk of $h_S$ into
        \emph{estimation error} (variance) $+$ \emph{approximation error} (bias) $+$ Bayes risk (statRisk, p.~4).
        Adding features lowers the bias but raises the variance; overfitting when variance dominates, underfitting when
        bias dominates (statRisk, p.~6). Near $d = m$ the variance diverges: this is the U-curve.
  \item \textbf{Interpolation threshold} ($d = m$): training error reaches $0$, test error peaks.
  \item \textbf{Over-parameterized regime} ($d > m$): training error is always $0$, the test error decreases again and
        can go below the classical minimum (paper, Fig.~1, p.~3). The paper's explanation (p.~3 and p.~6): beyond the threshold, richer classes contain interpolants with smaller norm, i.e.\ ``simpler'' ones (Occam's razor); the norm peaks at the threshold and then decreases (Fig.~2, p.~5). Classical theory, which predicts only the U, does not
        explain this part: this is the point of the paper.
\end{itemize}

\subsection{Planned plots}
\begin{enumerate}
  \item Training and test error against $d$ (minimum norm), vertical line at $d = m$, logarithmic $y$ axis
        (the peak spans orders of magnitude).
  \item On the synthetic data, a horizontal line at the Bayes risk $\sigma^2$ (Section 2) as a reference.
  \item $\norm{w}$ against $d$: it must peak at the same place (evidence for the mechanism of Section 5).
  \item Test error of minimum norm and ridge (several $\alpha$) on the same plot.
  \item The same plots for the real dataset, and a comparison between the two datasets.
\end{enumerate}

%--------------------------------------------------------------------
\section{Experimental design (preview)}

\paragraph{Relation to the paper.} The paper shows the phenomenon with random Fourier and random ReLU features (Sec.~2 and App.~C.1--C.2), neural networks (Sec.~2) and forests/boosting (Sec.~3), on MNIST and other classification datasets. It has no ``pure'' linear regression experiment, but the random features are trained by minimum-norm linear regression in feature space (p.~10): our experiment is its simplest version.

\subsection{Synthetic dataset: noisy linear labels}
\esterna{``weak features'' model: Belkin, Hsu, Xu, ``Two models of double descent for weak features'', 2019.}
\begin{itemize}
  \item $x \sim \mathcal{N}(0, I_D)$ with large $D$ (e.g.\ $D = 1000$), $y = w^{\star\top}x + \varepsilon$, Gaussian noise.
  \item The model with complexity $d$ uses only the first $d$ features. For $d < D$ the unused features act as extra
        noise (bias); for large $d$ the model has access to all the signal.
  \item Fixed $m$ (e.g.\ $m = 100$), $d$ from $1$ to $D$, denser grid around $d = m$.
  \item Mean (or median) over several repetitions with different seeds, with a variability band.
\end{itemize}

\subsection{Real dataset with random features}
\begin{itemize}
  \item A regression dataset with few features (e.g.\ California Housing, 8 features, also in the professor's repo).
  \item To push $d$ past $m$: random features $\phi(x) = \max(0, Wx)$ (ReLU) with a fixed Gaussian $W$, or random
        Fourier features as in the paper. The model stays \textbf{linear in the parameters} $w$.
        The course calls this idea \emph{feature expansion}: adding new features built from the original ones to reduce
        the bias of linear predictors (slides ``11. kernels'', p.~1).
  \item We subsample $m$ training examples (e.g.\ $m = 200$) and keep a separate test set.
\end{itemize}

\subsection{Rules against data leakage}
\begin{itemize}
  \item Train/test split \emph{before} any pre-processing.
  \item Mean and standard deviation for standardization computed \textbf{only} on the training set and applied to the test set.
  \item The random matrices $W$ do not depend on the data, so they create no leakage, but they must be generated with a fixed seed.
  \item Hyperparameters (e.g.\ $\alpha$) chosen only with training data.
\end{itemize}

%--------------------------------------------------------------------
\section{Likely oral questions, with short answers}

\begin{enumerate}
  \item \textbf{What is double descent?} As a function of complexity, the test error decreases, rises to a peak at the
        interpolation threshold ($d=m$), then decreases again in the over-parameterized regime.
  \item \textbf{Why is the peak exactly at $d = m$?} There $X$ is nearly square, $X^\top X$ has an almost-zero eigenvalue,
        the solution divides the noise by that eigenvalue and $\norm{w}$ blows up (slides 9, p.~5 and Section 5).
  \item \textbf{Why is the training error $0$ for $d \ge m$?} With $d \ge m$ and independent rows, $Xw = y$ has at least
        one exact solution: the model interpolates.
  \item \textbf{Which solution do you use for $d > m$ and why?} The minimum-norm $X^{+}y$: it is unique, interpolates,
        is the limit of ridge as $\alpha \to 0$ and the point GD converges to when started at $0$.
  \item \textbf{How did you compute $X^{+}$?} With \texttt{np.linalg.pinv}, based on the SVD. It is linear algebra, not a fit.
  \item \textbf{What does ridge do to the peak?} It replaces $\lambda_i$ with $\lambda_i + \alpha$ in the denominator:
        no noise amplification, peak damped or removed, at the cost of some bias.
  \item \textbf{How did you choose $\alpha$?} Fixed values compared side by side, and/or validation on the training set only. Never on the test set.
  \item \textbf{How do you vary $d$ on the real dataset?} With random features; the model stays linear in the parameters.
  \item \textbf{How do you avoid data leakage?} Split first, standardization statistics from the training set only,
        hyperparameters from the training set only.
  \item \textbf{Why does double descent contradict classical theory?} Classical theory predicts that beyond some point more
        complexity always hurts generalization (U-curve). In the $d > m$ regime the opposite happens.
  \item \textbf{Weak points of your work?} Linear models only; the real dataset needs random features, so the result depends
        on how they are built; few repetitions because of time constraints; arbitrary choice of $m$.
\end{enumerate}

%--------------------------------------------------------------------
\section{Solutions to the check questions}
\emph{Try to answer on your own first, then compare.}

\paragraph{Question 1.} Why can't the slides' formula $(X^\top X)^{-1}X^\top y$ be used for $d > m$, and what does \texttt{pinv(X) @ y} compute?

\smallskip
\noindent\textit{Answer.} $X^\top X$ is $d\times d$ but has rank at most $m < d$, so it is singular and not invertible
(equivalently, $m$ vectors cannot span $\R^d$, slides 9 p.~4). The system $Xw = y$ has infinitely many interpolating
solutions. \texttt{pinv(X) @ y} returns the one with minimum Euclidean norm, $X^\top(XX^\top)^{-1}y$: it is unique, lies in
the row space of $X$, and every other solution has a larger norm (Pythagoras). For $d \le m$ it coincides with the slides' formula.

\paragraph{Question 2.} Why does ridge have no peak at $d = m$?

\smallskip
\noindent\textit{Answer.} In the minimum-norm solution the components of $y$ are divided by the eigenvalues $\lambda_i$ of
$X^\top X$; at $d \approx m$ the smallest one is almost $0$ and the noise blows up. Ridge inverts $\alpha I + X^\top X$, whose
eigenvalues are $\lambda_i + \alpha \ge \alpha > 0$ (slides 9 p.~5): the denominator never drops below $\alpha$, so $\norm{w}$
and the test error stay bounded. The price is bias: with too large an $\alpha$, $w \to 0$ and the model underfits.

%--------------------------------------------------------------------
\section{Sources}
\begin{itemize}
  \item N.\ Cesa-Bianchi, notes ``Introduction'' (1.\ intro), p.~3--5; ``Statistical learning'' (4.\ statRisk), p.~1--6;
        ``Hyperparameter tuning and risk estimates'' (6.\ crossVal), p.~2--4; ``Kernels'' (11.\ kernels), p.~1.
  \item N.\ Cesa-Bianchi, slides ``Linear predictors'' (9.\ linear), version of 10/02/2026, p.~4--5.
  \item N.\ Cesa-Bianchi, slides ``Online Gradient Descent'' (10.\ online), version of 24/02/2026.
  \item M.\ Belkin, D.\ Hsu, S.\ Ma, S.\ Mandal, ``Reconciling modern machine learning practice and the classical
        bias-variance trade-off'', PNAS 2019 (arXiv:1812.11118). arXiv v2 (10/09/2019); page numbers refer to that PDF.
  \item M.\ Belkin, D.\ Hsu, J.\ Xu, ``Two models of double descent for weak features'', 2019 (arXiv:1903.07571). \emph{External.}
  \item T.\ Hastie, A.\ Montanari, S.\ Rosset, R.\ Tibshirani, ``Surprises in high-dimensional ridgeless least squares
        interpolation'', Annals of Statistics 2022. \emph{External.}
  \item P.\ Nakkiran, P.\ Venkat, S.\ Kakade, T.\ Ma, ``Optimal regularization can mitigate double descent'', ICLR 2021. \emph{External.}
  \item Assignment ``ML projects 25/26'', Assignment 4 (p.~8--9) and FAQ (p.~16).
\end{itemize}

\end{document}
